\section{Conclusion}
\label{sec:conclusion}

In 1941, Congress ended decades of apportionment disputes by permanently adopting the Huntington-Hill method. This mathematical formula resolved politically charged questions about seat allocation through transparent principles that no state could credibly challenge. Eighty years later, that formula continues to operate without controversy—not because it achieves perfect fairness, which mathematically cannot exist, but because it embodies procedural legitimacy through objectivity, transparency, and immunity to manipulation.

This paper demonstrates that similar principles can extend from apportionment to redistricting. If mathematical objectivity resolved \textit{how many} seats each state receives, similar approaches can determine \textit{where} district boundaries are drawn.

\subsection{Core Contributions}

We present a recursive graph bisection algorithm for congressional redistricting that operates exclusively on census data—population counts and geographic adjacency. The algorithm cannot access voter registration, election results, partisan affiliation, or any political information. Applied to all 50 states using 2020 Census data, it produces 435 congressional districts that:

\begin{itemize}
    \item Achieve strong population equality (mean deviation 2.79\%, 86\% of districts within $\pm$5\%)
    \item Guarantee geographic contiguity by construction (METIS contiguity constraint)
    \item Optimize compactness through edge-cut minimization
    \item Produce reproducible results from identical input data
    \item Possess structural immunity to partisan manipulation—the algorithm cannot see partisan data, therefore cannot intentionally gerrymander
\end{itemize}

Overlaying 2020 presidential election results reveals systematic Democratic advantage in district count (56.5\% Democratic-leaning districts versus 51.3\% national vote share). However, this pattern reflects well-documented geographic sorting—urban Democratic concentration and rural Republican dispersion—not algorithmic bias. The efficiency gap emerges from political geography and compactness optimization, not strategic manipulation.

\subsection{The Impossibility Defense}

Traditional anti-gerrymandering arguments focus on intent: did mapmakers deliberately manipulate boundaries for partisan advantage? The Supreme Court in \textit{Rucho v. Common Cause} found this standard non-justiciable, unable to distinguish legitimate political considerations from unconstitutional gerrymandering.

Algorithmic redistricting offers a stronger defense: \textit{impossibility}. The algorithm cannot gerrymander because it cannot see the information required for gerrymandering. Partisan manipulation requires knowledge of voter behavior—which precincts lean Democratic, which neighborhoods vote Republican. Our algorithm possesses none of this knowledge. It draws boundaries based on population and adjacency alone.

This structural blindness provides protection stronger than any institutional safeguard. Independent commissioners, no matter how well-intentioned, inevitably possess political knowledge that influences their decisions, consciously or unconsciously. An algorithm that has never seen political data cannot be influenced by it. If partisan patterns emerge in results, they reflect political geography and geometric optimization—not strategic design.

\subsection{Process Fairness Over Outcome Fairness}

Political scientists and legal scholars debate what constitutes "fair" political representation: proportional representation, partisan symmetry, competitive districts, or communities of interest. Different fairness metrics produce different—often contradictory—prescriptions.

We prioritize \textit{process fairness} over \textit{outcome fairness}. The algorithm may produce efficiency gaps, safe districts, or deviations from proportional ideals. These outcomes reflect political geography and compactness optimization, both of which constrain results regardless of redistricting method.

But the process itself is transparent, reproducible, uniform, and manipulation-resistant. These procedural virtues—mirroring Huntington-Hill—command legitimacy even when outcomes disappoint some normative frameworks.

Just as Huntington-Hill accepts that Wyoming receives 3× California's per-capita representation due to constitutional requirements, algorithmic redistricting accepts that geographic polarization creates efficiency gaps. The key difference between algorithmic redistricting and gerrymandering is not merely intent but \textit{impossibility}. The algorithm cannot favor Democrats or Republicans because it cannot distinguish them.

\subsection{Limitations and Extensions}

Algorithmic redistricting is not a panacea. Political geography constrains outcomes regardless of method. Communities of interest cannot be perfectly preserved. Voting Rights Act compliance requires explicit demographic constraints beyond pure geographic optimization.

Future enhancements could strengthen the approach:

\begin{itemize}
    \item \textbf{Edge-weighted optimization}: Minimizing total perimeter rather than edge counts could achieve 50--60\% compactness improvements, matching carefully-drawn manual maps.
    \item \textbf{Block-level implementation}: Using 8 million census blocks instead of 84,000 tracts would achieve population deviations approaching 0.1\%.
    \item \textbf{Multi-objective formulation}: Incorporating county boundaries, communities of interest, and other traditional criteria alongside compactness enables nuanced trade-off analysis.
    \item \textbf{VRA-constrained optimization}: Requiring specified numbers of majority-minority districts satisfies constitutional obligations while maintaining neutrality for other boundaries.
    \item \textbf{Ensemble methods}: Generating multiple plans enables statistical analysis of boundary robustness and outcome sensitivity to algorithmic choices.
\end{itemize}

These extensions are technically straightforward. The question is not algorithmic feasibility but normative prioritization: which criteria matter most, and how should trade-offs be balanced? A companion study~\citep{dellaLibera2026adaptive} examines adaptive tree selection methods that dynamically optimize the recursive bisection structure for improved fairness outcomes.

\subsection{Policy Implications}

States possess constitutional authority to adopt algorithmic redistricting without federal approval. Adoption pathways include ballot initiatives (direct voter approval), legislative enactment (voluntary delegation), judicial remedies (court-ordered alternatives to struck-down gerrymanders), and commission integration (algorithms implementing commissioner-specified criteria).

Public acceptance requires education: voters must understand that algorithmic redistricting prioritizes process fairness over outcome fairness. The algorithm may not produce proportional representation or competitive districts—geography constrains those outcomes—but it eliminates manipulation. Citizens can trust that boundaries reflect population and geography, not partisan calculation.

The Huntington-Hill precedent demonstrates feasibility. Before 1941, apportionment sparked recurring conflicts. After 1941, a mathematical formula commanded universal acceptance. The formula did not produce perfect fairness—which cannot exist—but it produced procedural legitimacy through objectivity and transparency.

Redistricting today faces similar legitimacy challenges. Public confidence in electoral fairness has eroded. Gerrymandering—real and perceived—undermines trust in representative democracy. Algorithmic approaches offer principled solutions grounded in mathematical objectivity rather than political negotiation.

\subsection{Broader Significance}

This work demonstrates that mathematical governance can extend beyond simple enumeration (apportionment) to complex geometric problems (boundary design). The principles underlying Huntington-Hill—objectivity, transparency, uniformity, reproducibility, and non-manipulability—transfer successfully to redistricting.

More broadly, algorithmic redistricting exemplifies how computational methods can address political problems resistant to traditional institutional solutions. When human judgment inevitably reflects political knowledge and strategic incentives, algorithms that operate without such knowledge provide structural protections impossible for human institutions to guarantee.

The approach is not anti-democratic. Democracy requires human deliberation about values, priorities, and trade-offs. But once democratic deliberation establishes criteria, algorithmic implementation ensures those criteria are faithfully applied without manipulation.

\subsection{Final Reflections}

The question is not whether algorithmic redistricting is perfect, but whether it is better than alternatives. On dimensions that matter most—objectivity, reproducibility, transparency, and elimination of direct partisan manipulation—the case is compelling.

Legislative redistricting has failed. Politicians drawing their own districts creates obvious conflicts of interest and has produced increasingly sophisticated gerrymanders as computational tools advanced. Public trust in electoral fairness has declined accordingly.

Independent commissions improve upon legislative redistricting but cannot guarantee neutrality. Commissioners possess political knowledge and make discretionary judgments. Outcomes depend on which individuals serve and what deliberations occur. Reproducibility and verification are impossible.

Algorithmic redistricting offers stronger guarantees. The algorithm cannot see political data, therefore cannot intentionally manipulate for partisan advantage. Anyone can re-run the calculation and verify results. Boundaries emerge from explicit mathematical principles rather than political negotiation.

Eighty years ago, Congress resolved apportionment disputes by delegating to a mathematical formula. The Huntington-Hill method succeeded not through perfection but through procedural legitimacy: transparent principles, uniform application, and structural immunity to manipulation.

Redistricting today requires similar solutions. Algorithmic approaches extend Huntington-Hill's philosophy from apportionment to boundary design, creating redistricting processes that command acceptance not through optimal outcomes but through defensible procedures.

If mathematical objectivity resolved how many seats each state receives, it can resolve where district boundaries are drawn. The technical tools exist. The precedent exists. The normative justification exists. What remains is political will to prioritize process fairness over strategic advantage—delegating not to algorithms blindly, but to mathematical principles that operate transparently, uniformly, and objectively in service of democratic legitimacy.
